% ! TEX root = ../m3-20-21.tex

\chapter{Serii Laurent și reziduuri}

\section{Serii de puteri. Serii Laurent}

Noțiunile privitoare la scrierea funcțiilor cu ajutorul seriilor de puteri,
întîlnite în cazul real prin serii Taylor, se pot generaliza și în cazul
complex. Situația este ceva mai problematică în acest caz, deoarece
argumentele sînt obiecte bidimensionale. În orice caz, multe dintre noțiunile
din cazul real pot fi preluate aproape fără modificare și în cazul complex.

\begin{definition}\label{def:serie-cpx}\index{serie!de puteri complexă}
  O serie de forma $ \ds\sum_{n \geq 0} a_n(z - z_0)^n $, cu $ z \in \CC $
  oarecare, $ z_0 \in \CC $ fixat și $ a_n \in \CC $ coeficienți, pentru
  orice $ n \in \NN $, se numește \emph{serie de puteri} centrată în $ z_0 $.
\end{definition}

Ca în cazul real, se poate calcula \emph{raza de convergență} a seriilor
de puteri cu una dintre formulele:
\[
  R = \dfrac{1}{\ds\lim_{n \to \infty} \sqrt[n]{|a_n|}} = %
  \lim_{n \to \infty} \left| \dfrac{a_n}{a_{n+1}} \right|.
\]

Atunci, deoarece lucrăm în planul complex, se definește \emph{discul de convergență}
al seriei prin:
\[
  B(z_0,R) = \{ z \in \CC \mid | z - z_0| < R \}.
\]

Știm, de asemenea, că:
\begin{itemize}
\item seria este absolut convergentă pe $ |z - z_0| < R $;
\item seria este divergentă pe $ |z - z_0| > R $;
\item seria este uniform convergentă pe $ |z - z_0| \leq \rho $, pentru
  orice $ \rho < R $,
\end{itemize}
iar pe frontiera discului trebuie testat separat.

În interiorul discului de convergență, suma seriei este o funcție olomorfă $ S(z) $
și au sens rezultate de forma derivării sau integrării termen cu termen.

\begin{definition}\label{def:analitica}\index{funcție!analitică}
  O funcție $ f : A \to \CC $ se numește \emph{analitică} dacă poate fi
  dezvoltată într-o serie de puteri complexă cu discul de convergență
  inclus în $ A $.
\end{definition}

Dacă acesta este cazul, avem formula cunoscută:
\[
  f(z) = \sum_{n \geq 0} \dfrac{f^{(n)}(z_0)}{n!} (z - z_0)^n.
\]

Are loc rezultatul important:
\begin{theorem}[Weierstrass-Riemann-Cauchy]\label{thm:wrc}
  O funcție complexă definită pe o mulțime deschisă este analitică dacă și
  numai dacă este olomorfă.
\end{theorem}

Deducem de aici că, dacă domeniul de definiție $ A $ al funcției complexe
studiate, $ f : A \to \CC $, este o mulțime deschisă, atunci olomorfia
(pe care o verificăm cu condițiile Cauchy-Riemann, de exemplu) este echivalentă
cu analiticitatea, verificată cu ajutorul seriilor de puteri. Rezultă
implicit că orice funcție olomorfă definită pe un deschis poate fi dezvoltată
în serie de puteri.

Dar pentru cazul complex, avem și serii ceva mai complicate:
\begin{definition}\label{def:laurent}\index{serie!Laurent}
  O serie de puteri de forma $ \ds\sum_{n \in \ZZ} a_n(z - z_0)^n $,
  cu $ a_n, z, z_0 \in \CC $ se numește \emph{serie Laurent} centrată în
  punctul $ z_0 \in \CC $.
\end{definition}

O serie Laurent este convergentă dacă și numai dacă atît partea pozitivă
$ a_n > 0 $, cît și partea negativă, $ a_n \leq 0 $ sînt simultan convergente.

Partea pozitivă se numește \emph{partea Taylor}, iar partea negativă se numește
\emph{partea principală}.

Pentru cazul Taylor, seriile uzuale pentru funcțiile elementare, împreună cu domeniile
de convergență, sînt date mai jos. În toate cazurile, $ x \in \RR $ se poate înlocui
cu $ z \in \CC $.
\begin{align*}
  e^x &= \sum_{n \geq 0} \dfrac{1}{n!} x^n, x \in \RR \\
  \dfrac{1}{1 - x} &= \sum_{n \geq 0} x^n, |x| < 1 \\
  \dfrac{1}{1 + x} &= \sum_{n \geq 0} (-1)^n x^n, |x| < 1 \\
  \cos x &= \sum_{n \geq 0} \dfrac{(-1)^n}{(2n)!} x^{2n}, x \in \RR \\
  \sin x &= \sum_{n \geq 0} \dfrac{(-1)^n}{(2n + 1)!} x^{2n + 1}, x \in \RR \\
  (1 + x)^\alpha &= \sum_{n \geq 0} \dfrac{\alpha(\alpha - 1)(\alpha - 2) \cdots (\alpha - n + 1)}{n!} x^n,%
                   |x| < 1, \alpha \in \RR \\
  \arctan x &= \sum_{n \geq 0} \dfrac{(-1)^n}{2n + 1} x^{2n + 1}, | x | \leq 1.
\end{align*}

Seriile pentru alte funcții se pot obține fie prin calcul direct, fie
prin substituții, fie prin derivare sau integrare termen cu termen a
seriilor de mai sus.

%%%%%%%%%%%%%%%%%%%%%%%%%%%%%%%%%%%%%%%%%%%%%%%%%%%%%%%%%%%%%%%%%%%%%%
\section{Singularități și reziduuri}

În cazul funcțiilor reale, domeniul de definiție trebuie ales atent astfel încît
să evităm \qq{punctele cu probleme}. Exemplele tipice sînt cele în care se
anulează numitorul unei fracții, cînd apar logaritmi sau radicali etc.
Pentru funcții reale, cel mult putem calcula asimptote verticale în acele
\qq{puncte cu probleme}. Dar în cazul complex, distincția se face mult mai fin.

\begin{definition}\label{def:sing-iz}\index{singularitate!izolată}
  Fie $ A \seq \CC $ o mulțime deschisă nevidă și $ f : A \to \CC $ o funcție
  complexă, olomorfă pe $ A $. Fie un punct $ z_0 \in \CC $.

  Punctul $ z_0 $ se numește \emph{punct singular izolat} (\emph{singularitate izolată})
  pentru $ f $ dacă există un disc centrat în $ z_0 $, de forma $ B(z_0, r) $, cu
  $ r \neq 0 $, astfel încît, dacă eliminăm centrul discului, funcția este
  olomorfă pe $ B(z_0, r) \setminus \{ z_0 \} $.
\end{definition}

Amintim că, dacă funcția este olomorfă pe \emph{discul punctat} (i.e.\ discul
$ B(z_0, r) $, din care am scos centrul), atunci ea are o dezvoltare în serie Laurent,
conform teoremei Weierstrass-Riemann-Cauchy.

\begin{definition}\label{def:sing-ap}\index{singularitate!aparentă}
  În condițiile și cu notațiile din definiția anterioară, punctul $ z_0 \in \CC $
  se numește \emph{punct singular izolat} (\emph{singularitate izolată}) dacă
  seria Laurent asociată funcției $ f $ are partea principală nulă, adică este
  o serie Taylor.
\end{definition}

\begin{definition}\label{def:pol}\index{singularitate!pol}
  În condițiile și cu notațiile din definiția anterioară, punctul $ z_0 \in \CC $
  se numește \emph{pol} dacă în seria Laurent asociată funcției $ f $ există
  un număr finit de termeni nenuli în partea principală. Indicele ultimului
  termen nenul $ m $ (în sensul că $ |m| $ este cel mai mare, iar $ m < 0 $)
  se numește \emph{ordinul polului}.
\end{definition}

\begin{definition}\label{def:sing-es}\index{singularitate!esențială}
  În condițiile și cu notațiile din definiția anterioară, punctul $ z_0 \in \CC $
  se numește \emph{punct singular esențial} (\emph{singularitate esențială})
  dacă în partea principală a seriei Laurent asociată funcției $ f $ există o
  infinitate de termeni nenuli.
\end{definition}

\vspace{0.3cm}

\textbf{Idee intuitivă:} Singularitățile, în general, sînt \qq{puncte cu probleme}.
Chiar și în cazul real există această noțiune, întîlnită în analiză. De exemplu,
un punct unghiular sau de întoarcere al unei funcții reale se numește singularitate.
Alt exemplu la îndemînă: găurile negre sînt interpretate în cosmologie (ramura fizicii
teoretice care utilizează geometria pentru studiul \qq{formei} Universului) ca
\emph{singularități ale spațiu-timpului}. Dar, din fericire, în majoritatea cazurilor,
singularitățile fie pot fi \qq{ignorate} sau \qq{eliminate} i.e. se poate extrage
informația de bază și din domeniul din care ele au fost scoase. Amintiți-vă, de exemplu,
că în analiza de liceu există o teoremă care afirmă că \emph{dacă se scoate un număr numărabil %
  de puncte din graficul unei funcții, integrala sa definită nu se schimbă}.

Similar este și cazul singularităților din analiza complexă: toate, cu excepția
celor esențiale, pot fi \qq{eliminate}, adică se poate extrage informație
foarte importantă despre comportamentul funcțiilor și în lipsa lor (sau, uneori,
chiar \emph{numai} din ele, ca în cazul reziduurilor, precum vom vedea).

\vspace{0.3cm}

Polii vor fi elementul central de studiu mai departe. Pentru ei, mai avem
o caracterizare utilă:

\begin{proposition}\label{pr:pol}
  În condițiile și cu notațiile de mai sus, punctul $ z_0 \in \CC $ este pol pentru
  funcția $ f $ dacă și numai dacă $ \ds\lim_{z \to z_0} f(z) = \infty $.
\end{proposition}

\textbf{Exemplu:} Putem vedea foarte simplu acest lucru pentru funcția
$ f(z) = \dfrac{z}{z - 1} $. Atunci $ z = 1 $ este pol simplu, deoarece
$ \ds\lim_{z \to 1} f(z) = \infty $ și mai mult, putem scrie dezvoltarea
în serie Laurent a funcției $ f $ în jurul lui $ z = 1 $, sub forma:
\[
  f(z) = \dfrac{1 + z - 1}{z - 1} = \dfrac{1}{z - 1} + 1,
\]
care are partea principală $ \dfrac{1}{z - 1} $, cu un singur termen
(care dă ordinul polului).

\begin{proposition}\label{pr:sing-es}
  În condițiile și cu notațiile de mai sus, punctul $ z_0 \in \CC $ este
  singularitate esențială dacă și numai dacă nu există
  $ \ds\lim_{z \to z_0} f(z) $.
\end{proposition}

Ajungem la o altă noțiune esențială:
\begin{definition}\index{singularitate!reziduu}
  În condițiile și cu notațiile de mai sus, fie $ z_0 $ o singularitate
  izolată a funcției $ f $. Coeficientul $ a_{-1} $ din seria Laurent
  a funcției $ f $ în coroana $ B(z_0; 0, r) $ se numește
  \emph{reziduul} funcției în $ z_0 $ și se notează $ \Rez(f,z_0) $.
\end{definition}

\begin{remark}\label{rk:calcul-rez}
  În multe situații, reziduurile se pot calcula simplu cu teorema \ref{thm:calc-rez},
  dar uneori sîntem obligați să folosim definiția, pentru că limitele complexe
  se calculează mult mai dificil decît în cazul real.
\end{remark}

Un exemplu este următorul: considerăm funcția $ f(z) = \dfrac{1}{z \sin z^2} $.
Vrem să calculăm $ \Rez(f, 0) $. Evident, acest punct este pol, deoarece
$ \ds\lim_{z \to 0} f(z) = \infty $, dar orice formulă am folosi din
teorema \ref{thm:calc-rez}, nedeterminarea nu este eliminată. Astfel că sîntem
obligați să folosim definiția, cu ajutorul seriei Laurent.

Ca în cazul seriei Taylor pentru funcția sinus, avem:
\begin{align*}
  \sin z &= \dfrac{z}{1!} - \dfrac{z^3}{3!} + \dfrac{z^5}{5!} - \dots \\
  \Rightarrow z \sin z^2 &= z \cdot \left( \dfrac{z^2}{1!} - \dfrac{z^6}{3!} + \dfrac{z^{10}}{5!} - \dots \right) \\
         &= z^3 \left( 1 - \dfrac{z^4}{3!} + \dfrac{z^8}{5!} - \dots \right).
\end{align*}

Atunci, inversînd, obținem funcția scrisă sub forma:
\[
  f(z) = z^{-3} \cdot \left( 1 - \dfrac{z^4}{3!} + \dfrac{z^8}{5!} - \dots \right)^{-1}.
\]

Vrem să inversăm seria din paranteză (admitem fără demonstrație că acest lucru
este posibil), pentru ca în final, să putem scrie toată funcția ca o serie de
puteri (în forma actuală nu este așa ceva!). Așadar, căutăm o serie de forma
$ \sum a_n z^n $, cu $ n \in \NN $, astfel încît:
\[
  \left( 1 - \dfrac{z^4}{3!} + \dfrac{z^8}{5!} - \dots \right) \cdot %
  \left(a_0 + a_1 z + a_2 z^2 + \dots \right) = 1.
\]
Prin identificarea coeficienților, rezultă $ a_0 = 1, a_2 = 0 $. Făcînd acum
înmulțirea, avem:
\[
  f(z) = \dfrac{1}{z^3} \sum_{n \geq 0} a_n z^n = %
  \dfrac{a_0}{z^3} + \dfrac{a_1}{z^2} + \dfrac{a_2}{z} + a_3 + \dots.
\]
Partea principală a seriei are 3 termeni nenuli, deci $ z = 0 $ este pol triplu
(intuitiv, simplu pentru $ z $ și dublu pentru $ \sin z^2 $), deci rezultă
$ \Rez(f, 0) = a_2 = 0 $.

%%%%%%%%%%%%%%%%%%%%%%%%%%%%%%%%%%%%%%%%%%%%%%%%%%%%%%%%%%%%%%%%%%%%%%
\section{Exerciții}

\indent\indent 1. Determinați dezvoltarea în serie Taylor sau Laurent pentru funcțiile următoare,
în jurul punctelor indicate $ z_0 $, precizînd și domeniile pe care sînt valabile:
\begin{enumerate}[(a)]
\item $ f(z) = \cos^2 z, z_0 = 0 $;
\item $ f(z) = \dfrac{z + 3}{z^2 - 8z + 15}, z_0 = 4 $;
\item $ f(z) = \sin \dfrac{1}{1 - z}, z_0 = 1 $;
\item $ f(z) = \dfrac{2z^2 + 9z + 5}{z^3 + z^2 - 8z - 12}, z_0 = 1 $;
\item $ f(z) = \dfrac{z - 1}{z - 2}, z_0 = 1 $;
\item $ f(z) = \dfrac{1}{z^2 - 3z + 2} $, pe domeniile $ |z| < 1 $,
  $ 1 < |z| < 2 $ și $ |z| > 2 $.
\end{enumerate}

\emph{Indicație:} Se pot folosi seriile uzuale, cu anumite substituții,
dar atenție la domeniile de convergență, precum și la eventualii poli ai
funcțiilor. A se vedea exemplul rezolvat de mai jos.

\vspace{0.3cm}

2. Determinați punctele singulare și precizați natura lor pentru funcțiile:
\begin{enumerate}[(a)]
\item $ f(z) = \dfrac{z^5}{(z^2 + 1)^2} $;
\item $ f(z) = z^3 \exp\left( -\dfrac{3}{z} \right) $;
\item $ f(z) = \sin \dfrac{1}{z} $;
\item $ f(z) = \dfrac{\sin 2z}{z^6} $;
\item $ f(z) = \dfrac{6z + 1}{z - 3} $;
\item $ f(z) = \dfrac{1}{z^2 \sin z} $;
\end{enumerate}

\emph{Indicație:} După identificarea \qq{punctelor cu probleme}, se cercetează
existență limitei către punctele respective și/sau se utilizează dezvoltarea
în serie Taylor/Laurent.

\vspace{0.3cm}

3. Calculați:
\begin{enumerate}[(a)]
\item $ \ds\int_{|z| = 1} \dfrac{z^2 + \exp(z)}{z^3} dz $;
\item $ \ds\int_\gamma (z + 1) \exp\left(\dfrac{1}{z+1}\right) dz $, unde
  $ \gamma: x^2 + y^2 + 2x = 0 $;
\item $ \ds\int_{|z| = r} \dfrac{\exp(z^{-1})}{1 - z} dz $, pentru $ r \in \left\{ \dfrac{1}{2}, 1, 2 \right\} $.
\end{enumerate}

\emph{Indicație:} Se folosește teorema reziduurilor.

\vspace{0.3cm}

\textbf{Exemplu rezolvat} pentru seria Laurent: Să se dezvolte funcția
\[
  f(z) = \dfrac{2z^2 + 3z - 1}{z^3 + z^2 - z - 1}
\]
în jurul originii și în jurul punctelor $ z = \pm 1 $.

\emph{Soluție:} Numitorul se descompune sub forma $ (z - 1)(z + 1)^2 $,
deci avem un pol simplu $ z = 1 $ și un pol dublu $ z = - 1 $.

Putem descompune funcția în fracții simple, sub forma:
\[
  f(z) = \dfrac{1}{z - 1} + \dfrac{1}{z + 1} + \dfrac{1}{(z + 1)^2}.
\]

Primii doi termeni sînt sume de serii geometrice, iar al treilea poate
fi obținut prin derivarea seriei geometrice:
\[
  \dfrac{1}{z + 1} = \sum (-1)^n z^n \Rightarrow %
  -\dfrac{1}{(z + 1)^2} = \sum (-1)^{n+1} (n + 1) z^n,
\]
după ce am făcut trecerea $ n \mapsto n + 1 $, întrucît seria derivatelor
ar fi pornit de la $ n = 1 $, dat fiind termenul $ z^{n-1} $.

Pe cazuri acum:

\textbf{Pentru $ |z| < 1 $:} funcția este olomorfă, deoarece nu avem
niciun pol în acest disc deschis. Atunci putem scrie direct seria ca sumă
a celor trei serii corespunzătoare:
\[
  f(z) = -\sum_{n \geq 0} z^n + \sum_{n \geq 0} (-1)^n z^n + %
  \sum_{n \geq 0} (-1)^n (n + 1)z^n = %
  \sum_{n \geq 0} z^n \left( -1 + (-1)^n + (-1)^n(n + 1)\right).
\]

\textbf{În jurul punctului $ z = - 1 $}, unde avem pol,
trebuie să rescriem funcția sub forma:
\[
  f(z) = \dfrac{1}{(z + 1)^2} + \dfrac{1}{z + 1} - %
  \dfrac{1}{2} \cdot \dfrac{1}{1 - \frac{z + 1}{2}},
\]
astfel punînd în evidență puteri (negative!) ale lui $ z + 1 $. Primii doi
termeni nu necesită prelucrare, iar pentru al treilea, întrucît
ne aflăm în domeniul de convergență a seriei geometrice, adică $ |z + 1| < 2 $
putem scrie:
\[
  \dfrac{1}{1 - \frac{z + 1}{2}} = \sum_{n \geq 0} \left( \dfrac{z + 1}{2} \right)^n.
\]

\textbf{În jurul punctului $ z = 1 $}, unde avem din nou pol, rescriem
funcția sub forma:
\[
  f(z) = \dfrac{1}{z - 1} + \dfrac{1}{2} \cdot \dfrac{1}{1 + \frac{z - 1}{2}} + %
  \dfrac{1}{4} \cdot \dfrac{1}{\left( 1 + \frac{z - 1}{2} \right)^2}.
\]

Primul termen nu necesită prelucrare, iar pentru al doilea, remarcăm că
ne aflăm în interiorul domeniului său de convergență, căci $ |z - 1| < 2 $
și atunci avem direct suma seriei geometrice:
\[
  \dfrac{1}{1 + \frac{z - 1}{2}} = \sum_{n \geq 0} (-1)^n \left( \dfrac{z - 1}{2} \right)^n.
\]

Al treilea termen se obține prin derivarea termen cu termen a seriei
de mai sus (ca în cazul de la început) și avem:
\[
  \dfrac{1}{\left(1 + \frac{z - 1}{2}\right)^2} = %
  \sum_{n \geq 0} (-1)^n \dfrac{(n + 1)(z - 1)^n}{2^n}
\]
și în fine, funcția se obține ca sumă a acestor trei serii, adică, după
calcule:
\[
  f(z) = \dfrac{1}{z - 1} + \sum_{n \geq 0} \dfrac{n + 3}{2^{n+2}} (z - 1)^n.
\]

%%% Local Variables:
%%% mode: latex
%%% TeX-master: "../m3-20-21"
%%% End:
